\documentclass{report}
\usepackage{pdfpages}
\usepackage{graphicx} 
\usepackage{dirtytalk}
\usepackage{lmodern}
\usepackage{booktabs}
\usepackage{listings}
\usepackage{float}
\usepackage{siunitx}
\usepackage{amsmath}
\usepackage{rotating}
\usepackage{pdflscape}
\usepackage{arydshln}
\usepackage[hidelinks]{hyperref}
\usepackage[ngerman]{babel}
\usepackage[a4paper, total={6in, 8in}]{geometry}
\usepackage{caption}
\begin{document}
\includepdf[pages=-]{TEP_Korrektur.pdf}
\chapter{Nachbesserung}
\section{Teilversuch I}
Aus der Steigung berechnet sich $x_0$ folgendermaßen:
\begin{equation}
    x_0^{\mathrm{steig}} = \frac{1}{m}
\end{equation}
Die Unsicherheit ist gegeben durch:
\begin{equation}
    \begin{split}
        \Delta x_0^{\mathrm{steig}} &= \left\lvert \frac{\partial x_0^{\mathrm{steig}}}{\partial m} \Delta m \right\rvert \\
                                    &= x_0^{\mathrm{steig}} \frac{\Delta m}{m}
    \end{split}
\end{equation}
Somit ergibt sich:
\begin{equation}
    \begin{split}
        x_0^{\mathrm{steig}} &= (74.074 \pm 0.658) \si{mm} \\
                                &= (74.1 \pm 0.7) \si{mm}
    \end{split}
\end{equation}
Damit gilt:
\begin{equation}
    x_0^{\mathrm{steig}} - 3\sigma = \qty{72.0}{mm}
\end{equation}
Der Wert ist also mit dem tatsächlichen Wert von $x_0 = \qty{72.0}{mm}$ gerade so im Rahmen einer $3\sigma$-Abweichung verträglich.
\end{document}